\documentclass[../../main]{subfiles}

\begin{document}

\section{Assinaturas Algébricas}
\label{sec:matematica:algebra:estruturas-algebricas:assinaturas}

Nesta seção introduzimos a noção de assinatura algébrica, que formaliza o tipo de operações presentes em uma estrutura algébrica, independentemente do conjunto subjacente.

\subsection{Símbolos de operação e aridade}

Um \textbf{símbolo de operação} é um símbolo abstrato associado a uma aridade fixa, isto é, um inteiro $n \in \mathbb{N}$, $n \geq 1$.

A aridade indica o número de argumentos da operação.

\begin{itemize}
	\item aridade $1$: operações unárias;
	\item aridade $2$: operações binárias;
	\item aridade $n$: operações n-árias.
\end{itemize}

\subsection{Definição de assinatura}

Uma \textbf{assinatura algébrica} é um conjunto
\[
	\Sigma = \{\sigma_i\}_{i \in I},
\]
onde cada símbolo $\sigma_i$ possui uma aridade fixa $n_i$.

Cada símbolo $\sigma_i$ será interpretado como uma operação:
\[
	\sigma_i^A : A^{n_i} \to A,
\]
quando aplicado a um conjunto $A$.

\subsection{Estruturas como interpretações de assinaturas}

Dada uma assinatura $\Sigma$, uma \textbf{estrutura algébrica de assinatura $\Sigma$} é um par
\[
	(A, \Sigma),
\]
onde:

\begin{itemize}
	\item $A$ é um conjunto não vazio;
	\item cada símbolo $\sigma \in \Sigma$ é interpretado como uma operação em $A$ com a aridade correspondente.
\end{itemize}

Ou seja, uma estrutura algébrica consiste em um conjunto juntamente com interpretações concretas dos símbolos de operação.

\subsection{Exemplos de assinaturas}

\begin{itemize}
	\item Grupos:
	      \[
		      \Sigma = \{\ast : 2\}
	      \]

	\item Anéis:
	      \[
		      \Sigma = \{+, \cdot : 2\}
	      \]

	\item Anéis com unidade:
	      \[
		      \Sigma = \{+, \cdot : 2, 0 : 0, 1 : 0\}
	      \]

	\item Espaços vetoriais (estrutura posterior):
	      envolve operações internas e externas:
	      \[
		      \Sigma = \{+, \cdot_{\text{ext}}\}
	      \]
\end{itemize}

\subsection{Observação estrutural}

A assinatura de uma estrutura algébrica determina apenas o tipo de operações envolvidas, mas não impõe propriedades sobre essas operações.

As propriedades (axiomas) serão introduzidas posteriormente como restrições adicionais sobre estruturas de uma mesma assinatura.

\subsection{Comentário conceitual}

A noção de assinatura permite separar dois níveis fundamentais da álgebra:

\begin{itemize}
	\item \textbf{nível sintático:} símbolos de operação e suas aridades;
	\item \textbf{nível semântico:} interpretações desses símbolos em um conjunto.
\end{itemize}

Essa separação será essencial para a definição de homomorfismos e para a comparação entre estruturas algébricas distintas.

\end{document}
